\begin{recipe}{Tarte tatin met wortels}
\kicker[Bijgerecht,Vegetarisch,Frans]{Bijgerecht · vegetarisch · Frans}
\index[register]{Tarte tatin met wortels}
\index[register]{Wortel}
\index[register]{Bladerdeeg}

\meta{1 uur \dvd Voor 4 personen}

\lettrine{E}{en} hartige tarte tatin van Rachel Khoo, met in boter gebakken wortels
onder een laag bladerdeeg. Een aardige afwisseling als bijgerecht, ook bij kerst.

\blockrule{Ingrediënten}
\begin{ingredients}
  \ing{10}{wortelen}\index[register]{Wortel}
  \ing{20 g}{ongezouten roomboter}
  \ing{1 el}{rodewijnazijn}
  \ing{½ el}{vloeibare honing}
  \ing{250 g}{vers bladerdeeg}\index[register]{Bladerdeeg}
  \ing{3 takjes}{verse tijm}
  \ing{1½ mespunt}{zout}
\end{ingredients}

\blockrule{Bereiding}
\tip{Lekker warm geserveerd met een groene salade.}
\begin{steps}
  \item Verwarm de oven voor op 180\,°C (boven- en onderwarmte).
  \item Schil de wortels als ze een dikke schil hebben, of was ze goed en dep droog.
        Smelt de boter in een grote koekenpan op laag vuur. Voeg zodra de boter
        begint te bruisen de wortels en tijm toe.
  \item Bak de wortels ca.\ 15 minuten en keer ze regelmatig zodat ze aan alle
        kanten kleuren. Haal de pan van het vuur en roer het zout, de honing en
        azijn erdoor. Schik de wortels op de bodem van een ronde taartvorm (18 cm).
  \item Rol het bladerdeeg uit tussen twee vellen bakpapier tot een dikte van
        5\,mm. Snijd een schijf uit die iets groter is dan de taartvorm. Leg het
        deeg op de wortels en stop de randen in. Snijd een klein kruis in het midden
        zodat stoom kan ontsnappen.
  \item Bak 30--35 minuten, tot het bladerdeeg gerezen en goudbruin is.
  \item Leg een groot bord op de vorm en draai de tarte tatin voorzichtig om.
        Serveer warm.
\end{steps}
\end{recipe}
